\documentclass[11pt]{article}

\usepackage{amsmath, amssymb, amsthm}
\usepackage{mathtools}
\usepackage{amsmath}
\usepackage{geometry}

\usepackage[T1]{fontenc}
\usepackage[polish]{babel}
\usepackage[utf8]{inputenc}

\geometry{margin=0.75in}
\title{\vspace*{-0.75in}\textbf{Lista 1}}
\author{Olaf Surgut \\ $345615$}

\begin{document}
\pagestyle{empty}
\maketitle
\thispagestyle{empty}

\section*{Zadanie 1}

Sprawdzić, że:
\begin{align*}
\text{(a)}\quad& \sum_{k = 0}^{n} \binom{n}{k} p^k (1-p)^{n-k} = (p + (1 - p))^n = 1 \\
\text{(b)}\quad& \sum_{k = 0}^{n} k \binom{n}{k} p^k (1-p)^{n-k} 
    = \sum_{k = 1}^{n} n \binom{n - 1}{k - 1} p^k (1-p)^{n-k}
    = n \sum_{k = 1}^{n} \binom{n - 1}{k - 1} p p^{k-1} (1-p)^{n-1-(k-1)} \\
    &= np \sum_{j = 0}^{n-1} \binom{n-1}{j} p^j (1-p)^{n-1-j} = np
\end{align*}


\section*{Zadanie 2}

Sprawdzić, że:
\begin{align*}
\text{(a)}\quad& \sum_{k=0}^{\infty} e^{-\lambda}\frac{\lambda^k}{k!} 
       = e^{-\lambda} \sum_{k=0}^{\infty} \frac{\lambda^k}{k!} 
      = e^{-\lambda} e^{\lambda} = 1 \\
\text{(b)}\quad& \sum_{k=0}^{\infty} k \cdot e^{-\lambda} \frac{\lambda^k}{k!}
      = \sum_{k=1}^{\infty} e^{-\lambda}\frac{\lambda^k}{(k-1)!}
      = \lambda \sum_{j=0}^{\infty} e^{-\lambda}\frac{\lambda^j}{j!}
      = \lambda
\end{align*}

\section*{Zadanie 3}

Funkcją $\Gamma$-Eulera nazywamy wartością całki:
\[
\Gamma(p) = \int_{0}^{\infty} t^{p-1}e^{-t}\, dt, \quad p > 0.
\]
Wykazać, że:
\[ \Gamma(p) = (p-1)\Gamma(p-1), \quad p \in \mathbb{R_{+}} \]

\[ \Gamma(p) = \int_{0}^{\infty} t^{p-1}e^{-t}\, dt \]
\[ 
\left\{
\begin{aligned}
u &= t^{p-1}, \qquad &dv&= e^{-t}\,dt\\ 
du &= (p-1)t^{p-2}\,dt, \quad &v& =-e^{-t} 
\end{aligned}
\right.
\]
\begin{align*}
\Gamma(p)
&= \Big[-t^{p-1}e^{-t}\Big]_{0}^{\infty} + (p-1)\int_{0}^{\infty} t^{p-2} e^{-t}\,dt 
= (p-1)\Gamma(p-1) 
\end{align*}
Podstawa indukcji dla $n = 0$.
\[
\Gamma(1)=\int_{0}^{\infty} t^{1-1}e^{-t}\, dt = \int_{0}^{\infty} e^{-t}\,dt = -e^{-t}\Big|_{0}^{\infty}=1=0!
\]

\section*{Zadanie 4}
\begin{align*}
(a)\;& \int_{0}^{\infty}\lambda e^{-\lambda x}\,dx= -e^{-\lambda x}\Big|_{0}^{\infty}=(0-(-1))=1 \\
(b)\;& \int_{0}^{\infty}x \lambda e^{-\lambda x}\,dx=
\end{align*}

\[\left\{
\begin{aligned}
u &= x, \quad &dv& = \lambda e^{-\lambda x}\,dx\\
du &= dx, \quad &v& = -e^{-\lambda x}
\end{aligned}
\right.\]

\begin{align*}
&= x (-e^{-\lambda x})\Big|_{0}^{\infty}-\int_{0}^{\infty} -e^{-\lambda x}\,dx=-\frac{1}{\lambda}e^{-\lambda x}\Big|_{0}^{\infty}=\frac{1}{\lambda}
\end{align*}

\section*{Zadanie 5}
Wykazać, że $D_n = n$, gdzie
\[
  D_n = 
\left|\,
\begin{matrix}
1      & -1 & -1 & \cdots & -1 \\
1      & \phantom{-}1 &    &        &    \\
1      &    &  \phantom{-}1 &        &    \\
\vdots &    &    & \ddots &    \\
1      &    &    &        & \phantom{-}1
\end{matrix}
\,\right| \overset{\text(1)}= 
\left|\,
\begin{matrix}
n      & -1 & -1 & \cdots & -1 \\
0      & \phantom{-}1 &    &        &    \\
0      &    &  \phantom{-}1 &        &    \\
\vdots &    &    & \ddots &    \\
0      &    &    &        & \phantom{-}1
\end{matrix}
\,\right| \overset{\text(2)}= 
\left|\,
\begin{matrix}
n      & 0 & 0 & \cdots & 0 \\
0      & 1 &   &        &    \\
0      &   & 1 &        &    \\
\vdots &   &   & \ddots &    \\
0      &   &   &        & 1
\end{matrix}
\,\right| = n \cdot \det(I_{n-1}) = n
\]

\begin{align*}
\text(1)&\quad c_{0}:=c_{0}-c_{1}-\ldots-c_{n}, \text{ gdzie $c_i$ to $i$-ta kolumna $D_n$}\\
\text(2)&\quad r_{0}:=r_{0}+r_{1}+\ldots+r_{n}, \text{ gdzie $r_i$ to $i$-ty wiersz $D_n$}
\end{align*}


\section*{Zadanie 6}
Niech
\[
I=\int_{-\infty}^{\infty} \exp\!\left(-\frac{x^{2}}{2}\right)\,dx.
\]
Wówczas
\[
I^{2}
=\left(\int_{-\infty}^{\infty} \exp\!\left(-\frac{x^{2}}{2}\right)\,dx\right)
 \left(\int_{-\infty}^{\infty} \exp\!\left(-\frac{y^{2}}{2}\right)\,dy\right)
=\int_{-\infty}^{\infty}\int_{-\infty}^{\infty}
\exp\!\left(-\frac{x^{2}+y^{2}}{2}\right)\,dy\,dx.
\]

Stosujemy podstawienie biegunowe
\[
x=r\cos\theta,\qquad y=r\sin\theta,
\]
gdzie \(r\in[0,\infty)\), \(\theta\in[0,2\pi)\), a wyznaczane pole wynosi
\[\text{Długość łuku $\times$ Grubość } = r\,d\theta\times\,dr\]
Zatem
\begin{align*}
I^{2}
&=\int_{0}^{2\pi}\int_{0}^{\infty}
\exp\!\left(-\frac{r^{2}}{2}\right)\, r \,dr\,d\theta \\
&=\left(\int_{0}^{2\pi} d\theta\right)
  \left(\int_{0}^{\infty} r\,e^{-r^{2}/2}\,dr\right).
\end{align*}
W całce po \(r\) podstawiamy \(u=\frac{r^{2}}{2}\), wtedy \(du=r\,dr\), więc
\[
\int_{0}^{\infty} r\,e^{-r^{2}/2}\,dr
=\int_{0}^{\infty} e^{-u}\,du
=1.
\]
Ostatecznie
\[
I^{2} = (2\pi)\cdot 1 = 2\pi.
\]

\section*{Zadanie 7}
Udowodnić, że:
\begin{align*}
\text{(a)}\quad& 
   \sum_{k=1}^{n}(x_k-\bar{x})^2
  =\sum_{k=1}^{n}x_k^2-2 x_k\bar{x}+\bar{x}^2
  =\sum_{k=1}^{n}(x_k^2)-2 n\bar{x}\bar{x}+n\bar{x}^2
  =\sum_{k=1}^{n}x_k^2-n\cdot\bar{x}^2  \\
\text{(b)}\quad& 
   \sum_{k=1}^{n}(x_k-\bar{x})(y_k-\bar{y})
  =\sum_{k=1}^{n}x_k y_k - x_k \bar{y} - \bar{x} y_k + \bar{x} \bar{y}
  =\sum_{k=1}^{n}(x_k y_k) - n \bar{x} \bar{y} - n \bar{x} \bar{y} + n \bar{x} \bar{y}
  =\sum_{k=1}^{n}x_k y_k - n\bar{x}\bar{y}
\end{align*}

\section*{Zadanie 8}
Dane są wektory $\vec{\mu}, X \in \mathbb{R}^n$ oraz macierz $\Sigma \in \mathbb{R}^{n\times n}$.
Niech $S = (X-\vec{\mu})^T\Sigma^{-1}(X-\vec{\mu})$ oraz $Y = A \cdot X$, gdzie macierz $A$ jest
odwracalna. Sprawdzić, że $S = (Y - A\vec{\mu})^T(A\Sigma A^T)^{-1}(Y - A\vec{\mu})$.

\begin{align*}
S &= (Y-A\vec{\mu})^T (A\Sigma A^T)^{-1} (Y-A\vec{\mu}) \\
  &= Y^T (A\Sigma A^T)^{-1} Y 
     - (A \vec{\mu})^T (A\Sigma A^T)^{-1} Y  \\
  &\quad - Y^T (A\Sigma A^T)^{-1} A \vec{\mu} 
     + (A \vec{\mu})^T (A\Sigma A^T)^{-1} A \vec{\mu} \\
  &= X^T \Sigma^{-1} X 
     - \vec{\mu}^T \Sigma^{-1} X 
     - X^T \Sigma^{-1} \vec{\mu} 
     + \vec{\mu}^T \Sigma^{-1} \vec{\mu} \\
  &= (X-\vec{\mu})^T \Sigma^{-1} (X-\vec{\mu}).
\end{align*}

\section*{Zadanie 9}
Udowodnić, że:
\[\int_{0}^{\infty}x^k \,\lambda \exp(-\lambda x)\,dx = \frac{k!}{\lambda^k}, \quad k = 0, 1, \ldots, \quad \lambda > 0\]


Podstawa indukcji $k = 0, 1$: (zadanie 4).

Indukcja $k \implies k + 1$:

\[\int_{0}^{\infty}x^k \,\lambda \exp(-\lambda x)\,dx = \]

\[\left\{
\begin{aligned}
u &= x^k, \quad &dv& = \lambda e^{-\lambda x}\,dx\\
du &= kx^{k-1}\,dx, \quad &v& = -e^{-\lambda x}
\end{aligned}
\right.
\]

\[
= x^k (-e^{-\lambda x}) \Big|_{0}^{\infty} - \int_{0}^{\infty} kx^{k-1} (-e^{-\lambda x})\, dx = \frac{k}{\lambda} \int_{0}^{\infty} x^{k-1} \lambda e^{-\lambda x}\,dx = \frac{k}{\lambda} \frac{(k-1)!}{\lambda^(k-1)} = \frac{k!}{\lambda^k} 
\]

\section*{Zadanie 10}
Obliczyć wartości całki
\[ \int_{0}^{\infty} x\,\exp(-\frac{x^2}{2})\,dx = \]
\[\left\{
\begin{aligned}
u = \frac{x^2}{2}, \quad du = x\,dx
\end{aligned}
\right.\]
\[ \int_{0}^{\infty} e^{-u}\,du = -e^{-u} \Big|_{0}^{\infty} = 0 - (-1) = 1\]

\end{document}